\section{Specializing a linked arithmetic seed on its given Hida ideal}

Assume global Artin reciprocity in the original normalization.
Fix the data in $\mathrm{SeedPtOfAlt}$: the algebraic closure
$\Omega/\mathbb Q$, $\ell$, $\iota$, $c,N,e,\psi$, the linked
$\eta,\sigma,g_1,B,M$ satisfying $\mathrm{LinkHyp}$, the CM field
$L$, frame and link $k$, and the concrete realization $R$ of
rank $n$ over $O=c_1.O$ with its given map of residue fields.
Write $F$ for the frame after this residue extension, $G$ for the
arithmetic group of the level $R.V$, and
\[
 G_L=\operatorname{Gal}(\Omega/L)
    \xrightarrow{\ p_V\ }G
    \xrightarrow{\ \pi\ }F.G_0,
 \qquad p_0=\pi p_V .
\]
The original group of the frame is unchanged by the residue
extension. Denote the twist character of the link on $G$ by
$\kappa$ and its pullback to $G_L$ by $k_a$, viewed in
$E=\overline{\mathbb Q}_\ell$. The given $\mathrm{IsKa}$ refers
to this $k_a$, the specified $\iota,\sigma,M$ and the same rank.

Set $\Psi=\mathrm{Chr.Psi}(\ell,e,\mathrm{dirG}(\psi))$ and
$\Theta=\epsilon^j\chi'$ on $G_{\mathbb Q}$, where
$\epsilon$ is the $\ell$-adic cyclotomic character. Write
$u_i=w(n,i)$, so $(u_i)$ is $(1,-1)$ or $(1,0,-1)$.
Assume $\ell\nmid\operatorname{ord}(\chi')$ and
$|\Theta(s)-\Psi(s)|_\ell<1$ for every $s$. Fix $m\ge2$
and an original $\mathrm{IsSeedAlt}$ representation $r$ on $G_L$
with residual characters $\Theta^{u_i}|_{G_L}$, away-from-$\ell$
inertia characters $(\chi')^{u_i}|_{G_L}$, inertia exponents
$m u_i$ above $\ell$, and scalar twist $k_a$. Write
$\widetilde r=k_a r$. The given automorphy field of the seed is the
original $\mathrm{IsAutCsdA}(\iota,\widetilde r)$.
Take any ideal $J$ of the actual determinant ring $B_d=R.B_0.B$
satisfying $R.B_0.\mathrm{PointTheoremBaseObject}(\iota,J)$.

% BEGIN P69A1
\paragraph{P69A1. Descent of the untwisted seed through its marked quotient.}
A prime-to-$\ell$ root of unity $z$ with $|z-1|_\ell<1$ equals
one: for its order $a$, the sum $1+z+\cdots+z^{a-1}$ has nonzero
residue $a$. At each valuation away from $\ell$, $\epsilon$ is
trivial on inertia. Thus on inertia where $\Psi(s)=1$ the
congruence of $\Theta$ gives $\chi'(s)=1$. The specified diagonal
action of $r$ on that inertia gives $r(s)=1$, including its
off-diagonal entries in the supplied diagonalizing basis.

Here this test gives descent from $G_L$ to precisely $F.G_0$.
Write $U=\mathrm{upl}(\Omega,L)$. The link identifies the image
of $G_L$ with $U\cap\ker\eta$, and its epimorphism
$U\to F.\Gamma$ identifies the sign quotient with $\eta$.
Consequently $p_0:G_L\to F.G_0$ is surjective. The kernel on
$U$ is the closed normal subgroup generated by the specified
away-from-$\ell$ inertia elements in $U$ with $\eta(s)=1$ and
$\bar\Psi(s)=1$. Every conjugate by $U$ again satisfies both
character conditions and belongs to inertia at the conjugate
valuation away from $\ell$. It therefore comes from $G_L$,
and the Teichm\"uller condition of the link upgrades
$\bar\Psi(s)=1$ to $\Psi(s)=1$ on this inertia. The preceding
test kills every such conjugate. The image of $\ker r$ in $U$
is closed by compactness, so it contains the required closed
normal subgroup. This proves the existence of $b:F.G_0\to
\mathrm{GL}_n(E)$ with $b p_0=r$. It is continuous because
$p_0$ is a continuous surjection of compact Hausdorff groups.
Irreducibility descends from $r$ along this surjection.

Take the supplied symmetric invertible matrix $P$ for the seed
pairing. Conjugation by $\sigma$ on $G_L$ induces the specified
conjugation on $F.G_0$ by the marked link. Lifting each $g$ proves
$b(g)^tP b(g^c)=P$. Set
$s(g)=\kappa(g)b(\pi g)$ on $G$. Then $s p_V=\widetilde r$,
and $s$ satisfies the original form equation with multiplier
$\epsilon^{1-n}$ because
$\kappa(g)\kappa(g^c)=\epsilon(g)^{1-n}$.
The map $p_V$ is surjective: lift a level-group element to the
absolute group and use the level's sign-kernel condition to
choose a lift in $G_L$. It is a continuous quotient map by
compactness. Since $r$ and $k_a$ are continuous, so are the
entries of $s$.
% END P69A1

% BEGIN P69A2
\paragraph{P69A2. The residual congruences and the two kinds of inertia.}
For $t\in G_L$ write $\Psi_L(t)$ for $\Psi$ of its image in
$G_{\mathbb Q}$, using the given inclusion $c.O\subseteq O$.
The $O$-valued characters
\[
             \beta_i(t)=\Psi_L(t)^{u_i}\kappa(p_Vt)
\]
reduce to the original labeled arithmetic characters
$R.r.\mathrm{chib}_i(p_Vt)$. Indeed $R.\mathrm{hf}$ matches the
two residue maps, $k.\mathrm{psi}$ identifies the framed residual
character, and the link's residual twist formula multiplies
its $u_i$th power by $\bar\kappa$. In $E$ their values are
$\Psi_L(t)^{u_i} k_a(t)$, all of norm one.

Write $x\equiv y$ when $|x-y|_\ell<1$. Congruent integral
elements have congruent sums and products. Congruent units also
have congruent inverses, by the formula for the difference
of two inverses. Hence these congruences allow every exponent
$u_i$, including $-1$. Multiplication of an $n$-dimensional
matrix by $k_a(t)$ multiplies the coefficient of $X^h$ in its
characteristic polynomial by $k_a(t)^{n-h}$ for $0\le h\le n$.
The seed's all-coefficient residual congruence and the congruence
$\Theta_L^{u_i}\equiv\Psi_L^{u_i}$ therefore give, coefficientwise,
\[
        \det(X-s(p_Vt))\equiv\prod_i(X-\beta_i(t)).
\]
The sign $(-1)^i$ in the $i$th determinant coefficient is the
same on both sides. Since $p_V$ is surjective, the congruence
holds against lifts of the prescribed residual coefficients at
\emph{every} element of $G$. Arbitrary two $O$-lifts of a fixed
residue differ by an element of norm less than one.

At each chosen place $v\mid\ell$, let $\alpha_{v,i}$ be the
ordered diagonal characters of the flag of $r$ furnished by
$\mathrm{OrdAt}$. On inertia they equal $\epsilon^{m u_i}$,
and on the entire decomposition group
$\alpha_{v,i}\equiv\Theta_L^{u_i}$. The diagonal characters for
the twisted flags are
$\gamma_{v,i}=\alpha_{v,i} k_a|_{D_v}$. Via the level local map
the same flag serves for $s$, because $s p_V=\widetilde r$.
It gives $\gamma_{v,i}\equiv\beta_i$ on $D_v$ in the prescribed
order. In particular each character and its inverse is integral.
The characters are continuous as diagonal entries in one fixed
flag basis. They kill the closure of the local commutator, so on
the specified $\mu_v$ they agree with the source's prescribed
character, which also equals one there by its commutator condition.

For an away-from-$\ell$ inertia element $t$ with
$R.r.\mathrm{chib}_i(p_Vt)=1$ for all $i$, the diagonal values
on $\widetilde r$ are $z_i=\chi'_L(t)^{u_i} k_a(t)$.
They are congruent to $\beta_i(t)\equiv1$ since $\epsilon(t)=1$.
Their common annihilating exponent is
$\operatorname{ord}(\chi')\,o$, where $o$ is the prime-to-$\ell$
exponent on away inertia from $\mathrm{IsKa}$. P69A1's
root-of-unity argument gives $z_i=1$, so $\widetilde r(t)=1$.
This kills the required $I_p(u)$ at every tame place: its
elements are images of inertia with all residual characters
one, and the marked tame set is away from $\ell$. It also
gives unramifiedness at the specific auxiliary place $v_0$.
Indeed $\ell\mid M$ and $v_0\nmid M$ imply $v_0\nmid\ell$,
and $R.v_0$ asserts residual triviality on its inertia for
all $i$. This latter conclusion will be used at $J$.
% END P69A2

% BEGIN P69A3
\paragraph{P69A3. A point over the original coefficient residue field.}
There exists an $O$-algebra map $x:B_d\to E$ realizing the full
determinant law of $s$ and all $\gamma_{v,i}$. We construct it
in the coefficient category of $B_d$. First choose a single finite
extension $D/\mathbb Q_\ell$ in $E$ containing the entries of all
$s(g)$. Indeed $E$ is the union of countably many finite
extensions (there are finitely many of each degree up to
isomorphism, and finitely many embeddings of each).
The closed subgroups where all entries lie in one such field
cover the compact domain $G$; one can take the fields closed
under finite compositum so each finite matrix is covered.
Baire's theorem gives one subgroup with interior and hence
finite index. Adjoining the matrix entries of finitely many
coset representatives then suffices. Enlarge $D$ by the
coefficient field of $O$ and the entries of the finitely many
flag bases and their inverses. All $\gamma_{v,i}$ then take
values in $D$.

Let $\mathfrak m_D$ be the maximal ideal of the integers $O_D$.
Use the subring $C=O+\mathfrak m_D\subset O_D$ as in HPT1.
It is finite over $O$, local with maximal ideal
$\mathfrak m_D$, and has exactly the residue field of $O$.
Completeness follows from finiteness and cofinality of the two
adic topologies, since $C/\mathfrak m_O C$ is Artinian local.
P69A2 puts every characteristic coefficient of every $s(g)$ in
$C$ with its specified residue, and puts every labeled
$\gamma_{v,i}$ and its inverse in $C$. These coefficient and
character functions are adically continuous: sufficiently small
differences lie in $\mathfrak m_D^a$ for each $a$, the same
powers as the maximal ideal powers in $C$.

For the \emph{law} use the integral word identity of QX5.
Each coefficient of the determinant of a generic finite linear
combination of group matrices is an integer polynomial in
characteristic coefficients of products of them. Each such
product is a group matrix. This descends the law to $C$,
and after reduction identifies it with the determinant of the
prescribed residual diagonal characters on all group-algebra
sums. Naturality and multiplicativity can be checked on generic
coefficients under the injection $C\hookrightarrow E$.
The form in P69A1 yields the polarized determinant relation.
The ordered ordinary product kills the full space: in the flag
basis its rightmost factor kills the last graded quotient,
the next factor kills the preceding graded quotient, and one
continues through the entire flag. This also proves the required
trace equation after any left group translate. The analogous
inertia trace relation holds since $s(g)=1$ on each prescribed
$I_p(u)$. Together with P69A2 these are all fields of the
original ordinary determinant pair over $C$. The universal
property gives a map $B_d\to C$; composing with $C\hookrightarrow
E$ gives $x$. Its pulled-back determinant on $G_L$ is that of
$\widetilde r$, and $x.\mathrm{psi}_{v,i}=\gamma_{v,i}$.
% END P69A3

% BEGIN P69A4
\paragraph{P69A4. Arithmetic weights and all forbidden cyclotomic ratios.}
Fix $v\mid\ell$ and use its chosen valuation together with the
chosen place $w$ for the local embeddings in the Hida weight.
Let $a(\tau_\infty)$ be the infinity exponents of the particular
algebraic Hecke character supplied for $k_a$ by $\mathrm{IsKa}$.
At each continuous $\tau:L_w\to E$ write
\[
 h_a(\tau)=-a\bigl((\iota\tau)|_L\bigr),\qquad
 h_{v,i}(\tau)=h_a(\tau)-m u_i .
\]
We check the locally algebraic assertion with the source convention
that the exponent in the local Artin product is the \emph{negative}
of the recorded weight. The smooth local component of the Hecke
character, pulled back from $L_w^\times$ to inertia and transported
by $\iota^{-1}$, has open kernel. Its inverse cancels that smooth
factor in the $\ell$-adic realization of $k_a$, leaving
$\prod_\tau\tau(a_w(s))^{-h_a(\tau)}$ for each compatible local
Artin map $a_w$ and inertia element $s$. The correction character
can be chosen once for all such maps. Given one map, their images
on inertia agree since they have the same global Artin image
and the local idele factor injects there (the local idele
injection from Chevalley's congruence theorem for units).
If the set of compatible maps is empty, the defining universal
condition holds with trivial correction. This gives precisely the
required smooth factor and weights for $k_a$.

In the same convention global Artin reciprocity gives cyclotomic
weight $-1$ at each embedding, with trivial smooth correction:
its value on inertia is the product of the local embeddings of the
Artin unit. Equivalently the $\ell$-adic realization of the global
absolute-value character is $\epsilon$, and its smooth local
absolute-value factor is one on units. By the seed hypothesis
$\alpha_{v,i}\epsilon^{-m u_i}$ is trivial on all inertia, hence
has weight zero there. Taking powers and multiplying the
identities proves that $\gamma_{v,i}$ is locally algebraic of
weight $h_{v,i}$ for the original chosen local datum. These
weights are strictly increasing in $i$ at every embedding, since
$u_i$ strictly decreases and $m\ge2$. Taking all $v$ gives the
Hida weight $h$ and the original $x.\mathrm{Arith}(h)$.

For \emph{any} $i,j$, equality
$\gamma_{v,i}=\gamma_{v,j}\epsilon$ would imply that
$\epsilon^{m(u_i-u_j)-1}=1$ on inertia. The exponent is nonzero
because $m\ge2$. Cyclotomic inertia at every such place of $L$
contains an element of infinite order in its cyclotomic image:
over $\mathbb Q_\ell$ the cyclotomic extension is totally ramified
and inertia maps onto $\mathbb Z_\ell^\times$. A positive power
of a lift of an infinite-order unit fixes the chosen number field
$L$ inside $\Omega$, since there are finitely many embeddings of
$L$. It remains in inertia and its value is still of infinite
order. This excludes all the ratios required in the point
condition, including $i=j$.
% END P69A4

% BEGIN P69A5
\paragraph{P69A5. Normalized realization and containment of the given ideal.}
Normalize $b$ of P69A1 on the frame without changing $x$.
Its local ordered characters come from $\gamma_{v,i}$ by
division by $\kappa$ at each local argument, then transport along
the link's local group isomorphism. The corresponding flags
and the vanishing of the ordered product follow by cancellation
of the unit scalar in each factor. The symmetric form $P$ from
P69A1 is already available. For clarity, the square root needed
for the marked-basis calculation of S68D4 exists here over
the same $C$ of P69A3. For $g\in G$ the framed trace at
$\pi(g)$ is $\kappa(g)^{-1}\operatorname{tr}(s(g))$,
which belongs to $C$ and reduces to
$\sum_i F.\Psi(\pi g)^{u_i}$ by P69A2. This trace is
star-invariant by the form equation for $b$.
At the marked $g_0$, put $z=F.\Psi(g_0)-F.\Psi(g_0)^{-1}$.
It is nonzero. Writing $t$ for the framed trace,
$t(g_0^2)-t(g_0g_0^*)$ reduces to $z^2$.
The source assumes odd $\ell$, so Hensel's lemma lifts the
square root with residue $z$. By the skew-adjoint calculation
of S68D4 the eigenvalues of $b(g_0)-b(g_0^*)$ are the
distinct $u_i$ times this root. Choose their eigenvectors as
the paired basis in S68D4. This makes the form antidiagonal
and the marked difference diagonal in the prescribed order.
In that basis $b$ and its local characters give an irreducible
normalized frame representation $b'$. Conjugacy preserves the
full determinant law, so
\[
 x_*D^{\mathrm{univ}}=\det\bigl(\kappa\cdot(b'\circ\pi)\bigr)
             \quad\text{on the arithmetic group }G .
\]

We may now use the given point condition of the specific $J$.
Its permitted rank interval $2\le n\le4$ contains both ranks
of $R$. P69A4 gives $x.\mathrm{Arith}(h)$, and P69A3 gives
$x.\mathrm{IsRep}(\widetilde r)$ on $G_L$ with all the specified
local characters. Continuity follows from the seed and
$\mathrm{IsKa}$, and the seed supplies the required automorphy
with the given $\iota$. P69A2 gives unramifiedness at $v_0$ and
all flags at $v\mid\ell$. P69A4 gives all the inequalities
against a cyclotomic translate. The conclusion of
$R.B_0.\mathrm{PointTheoremBaseObject}(\iota,J)$ is therefore
$J\subseteq\ker x$. The field $E$, the map $x$ and the
representation $b'$ furnish every field of
$R.T.\mathrm{IdealSpecializationPoint}(J)$. This proves the
native implication from global Artin reciprocity to
$\mathrm{SeedPtOfAlt}$.

For a Hida ideal the point condition is its $\mathrm{pt}$ field
as recorded in HIA1. With $\mathrm{Unr0}(J)$ and
$\mathrm{DimT}(J)$ this specialization now supplies the input
to S68D6 for that same $J$. The arithmetic inputs to this
implication include the given realization and $\mathrm{IsKa}$,
the congruent prime-to-$\ell$ character, and the
$\mathrm{IsSeedAlt}$ representation with its automorphic twist.
% END P69A5
